\section{方法}
\label{sec:methods}

\subsection{数据集与队列}

所有实验使用同一个经过整理的血吸虫病相关肝纤维化 B 型超声图像数据集,包含
来自 35 个中心、6,373 名患者的 108,709 张感兴趣区裁剪图像。研究方案经江苏省
血吸虫病防治研究所伦理审查委员会批准(批准号:JIPD-2019-017),所有参与者
在入组前均签署书面知情同意书。入组遵循前序队列研究的标准~\cite{xu2026sfibai}:
参与者招募自中国江苏省的血吸虫病流行县,覆盖完整严重度谱系,包括有确诊
日本血吸虫感染史者、登记在册的晚期血吸虫病患者以及未感染对照。性别信息未
采集,也未纳入研究设计。

每位患者按标准化六部位解剖协议采集,包括:剑突下矢状面(显示肝左叶)、左
肋缘下横切面(显示肝左叶及门静脉矢状部)、两个右肋缘下切面(分别显示第二
肝门三支肝静脉汇合处,以及带右肝静脉与膈肌的肝右叶)、两个右肋间斜切面
(分别显示带门静脉与胆囊的肝右叶,以及带右肾的肝右叶)。每个切面采集一至
三张图像,每张图像均带有归一化部位标签与纤维化评分。

纤维化严重程度在细粒度 36 级量表上标注,范围为 $0.0$ 至 $3.5$、步长 $0.1$。
四个临床分层源自中国血吸虫病相关肝纤维化的国家分级标
准~\cite{moh2000schistomanual,moh2006schistocriteria};其与国际超声指南
(主要针对\emph{曼氏血吸虫}感染制定)的对应关系,见 WHO 主持的巴塞尔超声
协议~\cite{richter2025basel}。用于临床报告时,连续量表在边界 $0.5$、$1.5$
与 $2.5$ 处映射到常规四级量表,边界值归入较高等级。图像标签取该图像标注
分数的最大值,患者标签取该患者所有图像的最大值,体现"患者严重程度以最重
发现为准"的临床惯例。

标注在 Labelme 中完成:围绕每个纤维化区域绘制矩形边界框,并同时赋予四级
分组与细粒度 36 级评分。为保证 0.1 级量表执行的一致性,标注医师接受了市级
与省级统一培训各一轮;标签随后由协调组集中复核,对不一致处给予纠正反馈,
部分病例还接受了盲法复审与二次评估。总体不一致率目标控制在 5\% 以下,存在
重大分歧的病例经共同复核达成共识分级。

图像数据质量控制剔除了:一批在非标准部位协议下采集的图像;感兴趣区裁剪或
标注框未通过校验的图像;经裁剪后内容哈希识别出的重复图像(包括所有跨患者、
跨中心或跨标签的冲突重复);无有效标注的图像(未被当作阴性处理);以及六
个标准部位中有效部位少于五个的患者。数据划分仅在清洗与去重之后生成,因此
没有任何患者或重复组跨越划分边界。

数据集在患者层面划分为:来自 4,906 名患者的 83,722 张训练图像、来自 1,227
名患者的 20,880 张验证图像,以及来自 240 名患者的 4,107 张测试图像。训练、
验证与测试集之间无患者重叠。各划分之间共享中心:测试集包含的全部四个中心
也出现在训练划分中。因此,测试集评估的是"已见中心内、未见患者"的泛化,
全文称其为患者级留出测试集。采用中心平衡的患者级指标,是为了降低中心样本
量差异的影响,而非估计在未见中心上的表现。

表~\ref{tab:cohort} 给出完整的逐划分构成。测试队列在 F0、F1、F2、F3 各
含 60 名患者,而训练与验证遵循队列分布,其中 F2 是患者层面最常见分级。由于
测试队列按分期平衡(每期 60 名患者),总体表现不应被解读为人群筛查场景下
的患病率加权估计。图像层面计数并不严格平衡,因为各患者贡献的图像数量不同。

\subsection{预处理}

所有图像统一缩放至 $512\times512$,并按 ImageNet 统计量归一化。主要配置
采用各向异性拉伸缩放;另设一个对照配置采用保持长宽比的信箱式(letterbox)
填充,从而将缩放约定的影响与建模贡献区分开来。

训练增广包括:以概率 $0.7$ 施加 $90^\circ$ 非零整数倍旋转;以概率 $0.5$
施加范围 $0.75$--$1.25$ 的饱和度抖动。不使用翻转或裁剪,因为在六部位采集
协议下水平翻转不具有解剖意义,而裁剪会破坏感兴趣区几何。成对比较的各配置
共享相同的样本顺序与逐样本增广抽样,使观察到的差异不来自数据呈现方式的
分歧。

对于病灶分支,弱监督目标在原始感兴趣区坐标系下构造为各标注框的并集掩码,
随后与图像同步缩放与旋转,使用最近邻插值以保持标签值。图像与掩码之间的
系统性错位被视为硬失败;单像素的边界离散化差异则予以容忍。

\subsection{SynAP-Fib 架构}

\subsubsection{分级骨干网络}

分级通路遵循 SFibAI 表述~\cite{xu2026sfibai}:以 \texttt{IMAGENET1K\_V2}
权重初始化的 ResNet-50 骨干网络,后接一个输出 36 个离散纤维化级别后验分布
的序数头。报告的连续预测值为后验期望
\begin{equation}
\hat{y} = \sum_{k=0}^{35} p_k \cdot \frac{k}{10},
\label{eq:score}
\end{equation}
由此得到与纤维化连续进展相一致的平滑估计,而非硬性级别指派。

SFibAI 基线在本研究的预处理、划分与评估协议下从头重训;因此,其在本文中的
表现不应与原始 SFibAI 实验设置下报告的数值直接比较。

\subsubsection{解剖先验分支}

SynAP-Fib 在该骨干网络上增设两个解剖先验分支,如
图~\ref{fig:architecture} 所示。

\emph{部位分支}是一个六类头,以交叉熵对归一化部位标签进行训练。真实部位
仅作为训练监督信号使用。推理时,分级通路读取的是\emph{预测的}部位后验,
从不读取部位元数据,因此该方法在可靠部位元数据不可用的场景中仍可部署。

\emph{病灶分支}预测一张空间注意力图,由弱框并集目标监督。仅纤维化分级为
阳性且至少有一个标注框的图像参与该目标;F0 图像与无框图像不计入病灶损失。

\subsubsection{门控残差注入}

两个分支通过同一机制影响分级。分支输出被转换为加权特征表示,以残差形式
加到分级表征上,其幅度由一个学习到的门控标量缩放:门控上界为
$\mathrm{gate\_max}=0.5$,初值为 $\mathrm{residual\_init\_scale}=0.001$,
使训练从接近基线通路的状态起步。联合配置将两个残差叠加。

分支与分级通路之间的梯度边界为双向隔离:分级损失的梯度不经注入回传到分支
头,分支损失的梯度也不经注入传播进分级通路。因此,各分支是解剖上下文的
来源,而非进入共享表征的额外梯度路径。所有门控超参数在每个配置中保持一致,
没有任何配置因门调控而获益。

\subsection{训练协议}

每个配置训练 120 个 epoch,不使用早停;优化器为 AdamW,学习率 $10^{-4}$,
权重衰减 $10^{-4}$,批大小 24,混合精度训练;学习率调度为每 15 个 epoch
乘以 $0.6$ 的阶梯调度。所有配置从相同的预训练初始化出发,各配置间共享的
模块在共同种子下逐张量地完全一致初始化。

分级目标为继承自 SFibAI 的混合序数损失,由邻域感知软标签项、期望回归项与
边界敏感惩罚项组成,权重分别为 $\alpha=1.0$、$\beta=0.02$ 与 $\gamma=0.02$,
软标签标准差为 $1.0$。部位目标与弱框目标各以 $0.1$ 的外层权重进入总损失。

所有模型在单张 NVIDIA GeForce RTX 5090 D GPU(CUDA 12.8)上训练,使用
PyTorch 2.11.0 与 torchvision 0.26.0(Python 3.10),启用自动混合精度。

\subsection{评估协议}

\subsubsection{临床目标风险}

评估采用复合临床目标风险(COR),同时惩罚连续分数误差与有临床后果的分级
错误。记绝对分数误差为 $e = |\hat{y}-y|$,分级距离为
$d = |\mathrm{grade}(\hat{y})-\mathrm{grade}(y)|$,则
\begin{align}
\mathrm{COR} ={}& 0.35\frac{e}{3.5} + 0.15\,\mathbb{I}(e>0.3) + 0.15\,\mathbb{I}(e>0.5) \nonumber\\
             & + 0.25\frac{d}{3} + 0.10\,\mathbb{I}(d\geq2).
\label{eq:cor}
\end{align}
所有样本等权,数值越低越好。

\subsubsection{聚合层级}

COR 在三个层级聚合。图像级风险 $R_{\mathrm{image}}$ 为逐图像 COR 的均值。
患者最大值风险 $R_{\mathrm{patient\text{-}max}}$ 在计算 COR 之前,对每位患者
的真实与预测图像分数对称地取最大值,与"患者严重程度以最重图像为准"的临床
惯例一致。中心平衡风险 $R_{\mathrm{cb}}$ 在每个中心内部计算患者最大值 COR,
再以与各中心患者数平方根成正比的权重合并各中心;其目的是防止大中心不成
比例地主导患者级总体评估,而非估计在未见中心上的表现。

预设的主要终点将三个层级合并为
\begin{equation}
R_{\mathrm{final}} = 0.4\,R_{\mathrm{image}} + 0.4\,R_{\mathrm{patient\text{-}max}}
                   + 0.2\,R_{\mathrm{cb}}.
\label{eq:rfinal}
\end{equation}
另保留一种对真实与预测分数对称计算的患者中位数聚合,作为稳健性分析。

\subsubsection{报告指标}

除 COR 外,我们报告平均绝对误差、$\pm0.3$ 与 $\pm0.5$ 容差内准确率、四分
级准确率与 macro-F1。另使用两个误差画像指标。超出 $0.5$ 容差的过量误差
指标 $\mathrm{TMAE} = \mathbb{E}[\max(e-0.5,\,0)]$ 在预测位于容差带内时为
零,仅随超出容差的误差部分增长。严重错误率在连续分数上定义为
$\Pr(e>1.0)$。判别与校准以 macro AUROC、macro AUPRC、Brier 分数与期望
校准误差概括。

对于病灶分支,我们报告阈值 $0.5$ 下相对于弱框并集的 Dice 与 IoU。这些指标
量化的是与弱边界框的一致程度,不可解读为病灶分割的真值表现,因为本队列不
存在像素级病灶标注。

\subsection{检查点选择与报告纪律}

检查点选择仅使用验证集。每个 epoch 都在验证集上完整评估;第 1--20 个
epoch 视为预热期,不参与选择。在第 21--120 个 epoch 中,选中的检查点是
验证 $R_{\mathrm{final}}$ 最低者;平局时依次取验证图像级 COR 更低者、epoch
更早者。不使用平滑、早停或选择器重采样。测试集在检查点、阈值或校准的选择
中不起任何作用。

本文报告的比较是 seed 2026 下的单种子比较,我们不将其作为运行间变异性的
证据。有两项指标由冻结的逐图像预测事后重算,并在全文如此标注:分期平均
绝对误差,以及 36 级上的二次加权 Cohen's $\kappa$。计算二次加权 Cohen's
$\kappa$ 前,连续的后验期望预测先映射到 36 个合法分级中最近的一个;未使用
后验 argmax 预测。
